\subsection{Conformal Set Aggregation (CSA)}
\label{sec2.4: CSA Method}
\noindent The standard split conformal framework in \ref{sec2.3.2: split conformal}, takes a scalar
non-conformity score $S(x, y) \in \mathbb{R}$. In this report, the prediction pipeline produces a score 
vector of scores $s \in [0, 1]^K$ for each phage, where $K$ is the number of functions. We may loose the 
geometric structure of the score distribution if we aggregate this vector to a scalar value before applying 
conformal prediction. For such settings, Ochoa Rivera, Pater and Tewari \cite{ochoarivera2024conformal} 
have proposed a score based aggregation (CSA) method for conformal prediction.

\subsubsection{Ensemble Score Vectors}
\label{sec2.4.1: ensemble score vectors}
\noindent Let $\mathcal{S} \subseteq [0, 1]^K$ be the score space and let $s = (s_1, \cdots, s_K) \in 
\mathcal{S}$ be the vector of scores produces by the $K$ models, for each phage. Now, the CSA method
treats the entire vector as the object of inference rather than reducing it into a single value. The main
idea is to define a region, called an envelope, $E \subseteq [0,1]^K$. This envelope consists of the region where
the score vectors of the true hosts for that particular phage live.

\vspace{0.15cm}
\noindent Mathematically, we can say that we apply a label dependent transformation $\phi_y: [0,1]^K \rightarrow
[0,1]^K$ to the score vector, before actually testing if it is part of the envelope. This transformation needs to be
such that the for a phage with the true label $y$, $\phi_y$ should be close to the origin in the score space. So,
we can write our prediction set as:
$$C(s) = \{y \in \mathcal{Y}: \phi_y(s) \in \hat t \dot E\},$$
\noindent where $\hat t$ is the global scaling factor that is calibrated on the calibration set \ref{sec2.3.2: split 
conformal}. The $\phi_y$ chosen in this work is defined in \todo{refer to the correct section}.


\subsubsection{The CSA Framework}
\label{sec2.4.2: csa framework}
\noindent The core innovation of the CSA method is in the calibration process. It splits the calibration data into
two subsets which have different roles. This is necessary for the exchangeability argument that provides us with the
coverage guarantee. If we used the same calibration points to construct both, we would get overfitting which would 
introduce bias which would invalidate the distruibution free coverage guarantee \cite{ochoarivera2024conformal}.

\vspace{0.15cm}
\noindent The first set is used to determine the geometric shape of the optimized envelope. Let's call it the 
'Shape Discover Set', it is used to map the distribution of the scores and then learn the geometry of the multi-dimensional
envelope $E \subseteq \mathbb{R}^K$. This envelope is such that the transformed score vectors of the true positives are
concentrated here. This set is only used to fit the shape of the conformal region without fixing the size.  

\vspace{0.15cm}
\noindent The second set or the 'Size Scaling Set' is used to determine the global scaling factor $\hat t \in 
\mathbb{R}^+$. This scaling factor either inflates or deflates the envelope, $E$ learned from the shape discovery set,
depending on the value of $\hat t$, so that we get the precise finite sample coverage guarantee.  This $\hat t$ is the 
smallest scaling factor such that the scaled envelope $\hat t \dot E$ achieves the required $1-\alpha$ coverage.
The scaled envelope is defined linearly as:
$$\hat t \cdot E = \{ t \cdot x : x \in E \}, \quad \hat t > 0$$

\noindent Finally, for a test input a candidate label, $y$ is included in the prediction set its transformed score vector
falls inside the scaled envelope:
$$C(s) = \{ y \in \mathcal{Y} : \phi_y(\mathbf{s}) \in \hat{t} \cdot E \}$$
where $\phi_y$ is a maps the raw scores into the score space.


\subsubsection{Limitations of CSA in Classification}
\label{sec2.4.3: csa limitations}
This CSA framework was originally developed and analyzed in a regression setting \cite{ochoarivera2024conformal}. 



\vspace{4cm}

While effective for multi-target regression, the original CSA formulation by Ochoa Rivera et al. has major limitations 
when applied to discrete classification problems:
 - Assumption of Convexity: The original paper assumes that $E$ is a convex hull or intersection of half-spaces. In 
                            bounded classification settings, the score vectors for distinct classes often concentrate in 
                            separate clusters near different boundaries of the $[0,1]^K$ hypercube. Forcing a convex 
                            envelope across these clusters creates a region that covers large areas of empty score space 
                            between them. This significantly inflates the volume of the envelope, leading to overly 
                            conservative, uninformative, and large prediction sets $|C(\mathbf{s})|$.
 - Iterative Calibration Inefficiency: The original method relies on iterative numerical optimization or binary search to 
                                       find the optimal scaling factor $\hat{t}$ over $S_2$. For high-dimensional score 
                                       spaces or massive metagenomic datasets, this iterative approach is computationally 
                                       expensive and lacks mathematical transparency.
 - Lack of Per-Class Guarantees: The original scaling calculation enforces only a marginal coverage guarantee, which can 
                                lead to poor predictive performance on highly imbalanced biological datasets.
These limitations motivate the core contributions of this thesis: designing structured non-convex envelopes, developing 
an exact analytical solution for $\hat{t}$, and implementing class-conditional adjustments to ensure robust statistical 
guarantees for phage-host range prediction.

